% !TEX root = ../../vzor.tex
%% Ukázka: citace, české uvozovky, tabulka
%% Vkládá ji soubor vzor/kapitoly/teorie.tex; ve vzor.pdf je pod nadpisem uvedeno „Zdroj ukázky".

\section{Citace}
\zdrojukazky{vzor/ukazky/citace.tex}
	Odkaz v závorkách: \parencite[viz][str. 900]{einstein}\\
	Odkaz: \cite{knuthwebsite}\\
	A odkaz pod čarou: \footcite[viz][str. 42]{latexcompanion}\\
	Odkaz na českou knihu: \cite{mareks}\\
	České uvozovky se sází příkazem \verb|\uv{text}|, například: \uv{příliš žluťoučký kůň úpěl ďábelské ódy}.\\
	Inline matematický výraz: $a^2 + b^2 = c^2$ nebo $\sum_{i=1}^{n} i = \frac{n(n+1)}{2}$.
	
	\begin{table}[htbp]
		\centering
		\caption{Naměřené hodnoty rychlosti}
		\label{tab:mereni}
		\begin{tabular}{lcc}
			\toprule
			Měření & Čas $t$ [\si{\second}] & Rychlost $v$ [\si{\meter\per\second}] \\
			\midrule
			1 & 1{,}25 & 3{,}84 \\
			2 & 1{,}30 & 3{,}69 \\
			3 & 1{,}22 & 3{,}93 \\
			\midrule
			Průměr & 1{,}26 & 3{,}82 \\
			\bottomrule
		\end{tabular}
	\end{table}

\lipsum[2]
